% ==============================================================================
% TRANG BÌA CHÍNH (TITLE PAGE)
% ==============================================================================
\begin{titlepage}
\begin{center}
    \textbf{\large BỘ GIÁO DỤC VÀ ĐÀO TẠO}\\[0.25cm]
    \textbf{\large TRƯỜNG ĐẠI HỌC PHENIKAA}\\[0.25cm]
    \textbf{\large KHOA KHOA HỌC MÁY TÍNH}\\[0.8cm]

    \vspace{0.8cm} % Chỗ chèn Logo Trường Đại học Phenikaa (nếu có: \includegraphics[width=0.25\textwidth]{figures/logo.png})

    \textbf{\Large BÁO CÁO ĐỒ ÁN TỐT NGHIỆP ĐẠI HỌC}\\[0.3cm]
    \textbf{\large NGÀNH: KHOA HỌC MÁY TÍNH}\\[0.2cm]
    \textbf{\large CHUYÊN NGÀNH: TRÍ TUỆ NHÂN TẠO}\\[0.8cm]

    \begin{onehalfspace}
        {\bfseries\large NGHIÊN CỨU VÀ XÂY DỰNG TRÒ CHƠI TRINH THÁM GIẢI ĐỐ TRÊN NỀN TẢNG WEB, ỨNG DỤNG MÔ HÌNH NGÔN NGỮ LỚN (LLM) ĐỂ SINH TỰ ĐỘNG CÁC BỘ HỒ SƠ VỤ ÁN CÓ TÍNH LOGIC VÀ NHẤT QUÁN}\\[0.35cm]
        {\normalsize\textbf{Hướng tiếp cận Neuro-Symbolic kết hợp Hệ đa tác tử và Bộ kiểm định hình thức}}\\[0.25cm]
        {\small\textit{(Research and Development of a Web-Based Detective Puzzle Game Using Large Language Models for Procedural Generation of Logical and Consistent Case Dossiers: A Neuro-Symbolic Multi-Agent Approach)}}
    \end{onehalfspace}

    \vfill

    \begin{table}[h]
        \centering
        \begin{tabular}{ll}
            \textbf{Sinh viên thực hiện:} & \textbf{Nguyễn Ngọc Công Anh} \\
            \textbf{Mã số sinh viên:}     & \textbf{22010201} \\
            \textbf{Lớp:}                & \textbf{K16-KHMT} \\
            \textbf{Cán bộ hướng dẫn:}    & \textbf{TS. Nguyễn Văn Tánh}
        \end{tabular}
    \end{table}

    \vspace{0.4cm}
    \textbf{\large HÀ NỘI – NĂM 2026}
\end{center}
\end{titlepage}

% ==============================================================================
% LỜI CAM ĐOAN
% ==============================================================================
\newpage
\chapter*{LỜI CAM ĐOAN}
\addcontentsline{toc}{chapter}{LỜI CAM ĐOAN}

Tôi xin cam đoan bản báo cáo đồ án tốt nghiệp với đề tài \textit{``Nghiên cứu và xây dựng trò chơi trinh thám giải đố trên nền tảng Web, ứng dụng mô hình ngôn ngữ lớn (LLM) để sinh tự động các bộ hồ sơ vụ án có tính logic và nhất quán''} là công trình nghiên cứu độc lập của riêng tôi dưới sự định hướng và hướng dẫn khoa học trực tiếp của TS. Nguyễn Văn Tánh.

Toàn bộ các số liệu đo đạc thực nghiệm, kết quả thống kê benchmark và mã nguồn hệ thống được trình bày trong luận văn này là trung thực, phản ánh đúng kết quả thu được trong quá trình nghiên cứu và cài đặt. Mọi sự kế thừa từ các công trình nghiên cứu trước đây, các thư viện mã nguồn mở, cũng như các tài liệu tham khảo khoa học đều được trích dẫn và ghi nhận tường minh theo đúng chuẩn mực liêm chính học thuật quốc tế. Tôi xin chịu hoàn toàn trách nhiệm trước Hội đồng chấm tốt nghiệp và Nhà trường về tính nguyên bản của công trình này.

\vspace{1.5cm}
\begin{flushright}
    \textit{Hà Nội, ngày 17 tháng 11 năm 2026}\\[0.2cm]
    \textbf{Sinh viên thực hiện}\\[1.5cm]
    \textbf{Nguyễn Ngọc Công Anh}
\end{flushright}

% ==============================================================================
% LỜI CẢM ƠN
% ==============================================================================
\newpage
\chapter*{LỜI CẢM ƠN}
\addcontentsline{toc}{chapter}{LỜI CẢM ƠN}

Để có thể hoàn thành chương trình đào tạo ngành Khoa học máy tính tại Trường Đại học Phenikaa và hoàn thiện bản báo cáo đồ án tốt nghiệp này, tôi đã nhận được sự dạy bảo, hỗ trợ và động viên to lớn từ các Thầy Cô giáo, gia đình và bạn bè.

Lời đầu tiên, tôi xin bày tỏ lòng biết ơn sâu sắc và chân thành nhất tới \textbf{TS. Nguyễn Văn Tánh}, người Thầy đã tận tâm hướng dẫn, truyền thụ phương pháp tư duy khoa học, định hướng kiến trúc hệ thống và luôn đưa ra những câu hỏi phản biện nghiêm khắc giúp tôi trưởng thành đáng kể cả về năng lực học thuật lẫn tư duy kỹ thuật trong suốt quá trình thực hiện đồ án.

Tôi xin chân thành cảm ơn Ban Giám hiệu Trường Đại học Phenikaa, Ban Chủ nhiệm Khoa Khoa học máy tính cùng toàn thể quý Thầy Cô giáo đã truyền thụ những nền tảng tri thức quý báu về Toán học, Cấu trúc dữ liệu và giải thuật, Học máy, Xử lý ngôn ngữ tự nhiên và Kỹ nghệ phần mềm trong suốt bốn năm học tập dưới mái trường.

Con xin bày tỏ lòng biết ơn sâu nặng tới Bố Mẹ và gia đình, những người luôn là điểm tựa tinh thần vững chắc, tạo mọi điều kiện tốt nhất để con yên tâm học tập và theo đuổi định hướng nghiên cứu khoa học.

Cuối cùng, tôi xin gửi lời cảm ơn tới những người bạn đồng môn đã luôn đồng hành, chia sẻ kinh nghiệm, hỗ trợ kiểm thử hệ thống và thảo luận chuyên môn trong suốt thời gian qua.

\vspace{1.5cm}
\begin{flushright}
    \textit{Hà Nội, tháng 11 năm 2026}\\[0.2cm]
    \textbf{Sinh viên}\\[1.5cm]
    \textbf{Nguyễn Ngọc Công Anh}
\end{flushright}

% ==============================================================================
% TÓM TẮT ĐỒ ÁN (TIẾNG VIỆT VÀ TIẾNG ANH)
% ==============================================================================
\newpage
\chapter*{TÓM TẮT ĐỒ ÁN}
\addcontentsline{toc}{chapter}{TÓM TẮT ĐỒ ÁN}

Sự phát triển nhanh chóng của các mô hình ngôn ngữ lớn (LLM) mở ra tiềm năng đáng kể cho lĩnh vực sinh nội dung thủ tục (PCG) trong trò chơi điện tử. Tuy nhiên, việc ứng dụng LLM nguyên khối vào thể loại trò chơi giải đố trinh thám suy diễn gặp phải nhiều rào cản bản chất như hiện tượng ảo giác nhân quả, mâu thuẫn mốc thời gian tử vong sinh học và vi phạm nguyên tắc trinh thám công bằng theo mười điều răn của Ronald Knox. Cơ chế sinh tự quy hồi theo xác suất của LLM không đủ để bảo đảm tính nhất quán của chuỗi suy luận hình thức nhiều bước.

Đồ án này đề xuất một khung kiến trúc lai Neuro-Symbolic kết hợp chuỗi 4 tác tử tuần tự (\texttt{OutlineAgent}, \texttt{CharacterAgent}, \texttt{ClueAgent}, \texttt{DeductionAgent}) với bộ kiểm định quy tắc tất định. Hệ thống tách biệt rõ ràng giữa dữ liệu chân lý gốc với dữ liệu trình bày cho người chơi, tích hợp phương trình làm lạnh thi thể hai hàm mũ của Henssge (giải số bằng phương pháp chia đôi) để khóa cứng cửa sổ thời gian tử vong (TOD), đồng thời sử dụng biểu diễn đồ thị hai phía để kiểm tra tính duy nhất của thủ phạm và sự tồn tại của manh mối bẻ gãy ngoại phạm. Khi phát hiện sai phạm, cơ chế sửa lỗi cục bộ định tuyến tín hiệu chẩn đoán tới đúng tác tử liên quan với độ phức tạp ngữ cảnh hằng số $\mathcal{O}(1)$ nhằm tránh tái sinh toàn bộ hồ sơ ($\mathcal{O}(N)$). Hệ thống vận hành trên kiến trúc dịch vụ Web RESTful API (\texttt{FastAPI}, giao thức HTTP/JSON, bảo mật HMAC-SHA256) kết hợp lưu trữ SQLite (SQLAlchemy 2.0) và giao diện điều tra đa nền tảng trên Unity 6 kết xuất bằng đường ống kết xuất phổ quát (URP) hỗ trợ chuẩn WebGL trình duyệt và máy trạm.

Kết quả thực nghiệm qua 30 phiên thử nghiệm độc lập trên 3 dòng mô hình (\texttt{gh/gpt-4o-mini}, \texttt{gh/gpt-4o}, \texttt{ollama/gemma4:31b-cloud}) cho thấy toàn bộ 30 kịch bản đạt \textbf{100\% tỷ lệ hợp lệ cú pháp JSON} dưới chuẩn Pydantic v2, nâng tỷ lệ vượt qua kiểm định logic lên mức \textbf{90\% -- 100\%}, chặn lọc được lỗi tham chiếu thực thể phát sinh ở \texttt{gh/gpt-4o}, và giữ độ trễ trung vị P50 ở mức \textbf{10.73s -- 10.88s} đối với hai mô hình phản hồi nhanh (\textbf{19.53s} đối với \texttt{gh/gpt-4o}).

\vspace{0.5cm}
\noindent\textbf{Từ khóa:} \textit{Mô hình ngôn ngữ lớn, Hệ đa tác tử, Neuro-Symbolic AI, Sinh nội dung thủ tục, Trò chơi trinh thám, Mô hình nhiệt Henssge, Đồ thị hai phía, Pydantic v2, Unity URP.}

\newpage
\chapter*{ABSTRACT}
\addcontentsline{toc}{chapter}{ABSTRACT}

The rapid evolution of Large Language Models (LLMs) presents substantial opportunities for Procedural Content Generation (PCG) in interactive video games. However, applying monolithic LLMs to deductive mystery games encounters fundamental obstacles, including causal hallucinations, biological post-mortem interval contradictions, and frequent violations of classic ``Fair-Play Mystery'' conventions (Knox's Decalogue). The probabilistic autoregressive nature of LLMs cannot guarantee strict formal consistency across multi-hop deductive inference chains on its own.

This thesis presents a \textbf{Neuro-Symbolic Hybrid Framework} that combines a 4-stage sequential Multi-Agent pipeline (\texttt{OutlineAgent}, \texttt{CharacterAgent}, \texttt{ClueAgent}, \texttt{DeductionAgent}) with a deterministic symbolic rule-based validator engine. The architecture decouples the underlying ground-truth state from player-facing dossier presentation and verification. It incorporates Henssge's two-exponential forensic cooling model solved numerically via the bisection method to enforce post-mortem Time of Death ($TOD$) windows as physical invariants, while using bipartite graph modeling to verify culprit uniqueness and alibi-breaking clue reachability. When constraint violations are detected, a localized patch repair mechanism isolates the responsible component with $\mathcal{O}(1)$ repair-context complexity rather than regenerating the entire dossier ($\mathcal{O}(N)$). The complete system operates on a distributed Web service architecture using a Python FastAPI RESTful API (HTTP/JSON, HMAC-SHA256) backed by SQLite (SQLAlchemy 2.0 WAL) and a cross-platform Unity 6 URP 2D corkboard client targeting both browser WebGL and standalone runtime environments.

Empirical evaluation across 30 independent trials benchmarking three representative models (\texttt{gh/gpt-4o-mini}, \texttt{gh/gpt-4o}, \texttt{ollama/gemma4:31b-cloud}) achieves \textbf{100.0\% syntactic schema validity} under Pydantic v2, maintains \textbf{90.0\% -- 100.0\% logic validation pass rates} while reliably flagging referential hallucination errors before client delivery, and records median P50 end-to-end generation latencies of \textbf{10.73s -- 10.88s} on fast open-weights and compact commercial models (\textbf{19.53s} on \texttt{gh/gpt-4o}).

\vspace{0.5cm}
\noindent\textbf{Keywords:} \textit{Large Language Models, Multi-Agent Systems, Neuro-Symbolic AI, Procedural Content Generation, Detective Mystery Games, Henssge Cooling Model, Bipartite Graph Solvability, Pydantic v2, Unity URP.}

% ==============================================================================
% MỤC LỤC VÀ DANH MỤC DANH SÁCH
% ==============================================================================
\newpage
\tableofcontents

\newpage
\listoftables
\addcontentsline{toc}{chapter}{DANH MỤC BẢNG BIỂU}

\newpage
\listoffigures
\addcontentsline{toc}{chapter}{DANH MỤC HÌNH VẼ}

% ==============================================================================
% DANH MỤC TỪ VIẾT TẮT
% ==============================================================================
\newpage
\chapter*{DANH MỤC CÁC TỪ VIẾT TẮT}
\addcontentsline{toc}{chapter}{DANH MỤC CÁC TỪ VIẾT TẮT}

{\small
\begin{longtable}{c l p{5.5cm} p{5.8cm}}
\toprule
\textbf{STT} & \textbf{Ký hiệu} & \textbf{Thuật ngữ tiếng Anh} & \textbf{Giải nghĩa tiếng Việt} \\
\midrule
\endfirsthead
\toprule
\textbf{STT} & \textbf{Ký hiệu} & \textbf{Thuật ngữ tiếng Anh} & \textbf{Giải nghĩa tiếng Việt} \\
\midrule
\endhead
\bottomrule
\endfoot
1  & \textbf{ACID}    & Atomicity, Consistency, Isolation, Durability & Tính nguyên tử, nhất quán, cô lập, bền vững \\
2  & \textbf{ADR}     & Architecture Decision Record                  & Bản ghi quyết định kiến trúc \\
3  & \textbf{AI}      & Artificial Intelligence                       & Trí tuệ nhân tạo \\
4  & \textbf{API}     & Application Programming Interface             & Giao diện lập trình ứng dụng \\
5  & \textbf{CoT}     & Chain-of-Thought                              & Kỹ thuật suy luận theo chuỗi \\
6  & \textbf{FSM}     & Finite State Machine                          & Máy trạng thái hữu hạn \\
7  & \textbf{GPT}     & Generative Pre-trained Transformer            & Mô hình Transformer tạo sinh tiền huấn luyện \\
8  & \textbf{HMAC}    & Hash-based Message Authentication Code        & Mã xác thực thông điệp dựa trên hàm băm \\
9  & \textbf{HTTP}    & Hypertext Transfer Protocol                   & Giao thức truyền tải siêu văn bản \\
10 & \textbf{JSON}    & JavaScript Object Notation                    & Định dạng trao đổi dữ liệu văn bản có cấu trúc \\
11 & \textbf{LAN}     & Local Area Network                            & Mạng máy tính cục bộ \\
12 & \textbf{LLM}     & Large Language Model                          & Mô hình ngôn ngữ lớn \\
13 & \textbf{NPC}     & Non-Player Character                          & Nhân vật không phải người chơi \\
14 & \textbf{P50}     & 50th Percentile (Median)                      & Phân vị thứ 50 (Trung vị) \\
15 & \textbf{P95}     & 95th Percentile                               & Phân vị thứ 95 \\
16 & \textbf{PCG}     & Procedural Content Generation                 & Sinh nội dung theo thủ tục \\
17 & \textbf{PMI}     & Post-Mortem Interval                          & Khoảng thời gian sau tử vong \\
18 & \textbf{RAM}     & Random Access Memory                          & Bộ nhớ truy cập ngẫu nhiên \\
19 & \textbf{REST}    & Representational State Transfer               & Kiến trúc truyền trạng thái đại diện \\
20 & \textbf{SHA-256} & Secure Hash Algorithm 256-bit                 & Giải thuật băm an toàn 256-bit \\
21 & \textbf{TOD}     & Time of Death                                 & Cửa sổ thời gian tử vong \\
22 & \textbf{URP}     & Universal Render Pipeline                     & Đường ống kết xuất phổ quát trong Unity \\
23 & \textbf{WAL}     & Write-Ahead Logging                           & Cơ chế ghi nhật ký trước trong cơ sở dữ liệu \\
\end{longtable}
}
